% Generated draft: requires mathematical and source review.
\documentclass[11pt]{article}
\usepackage[margin=25mm]{geometry}
\usepackage{fontspec}
\setmainfont{Latin Modern Roman}
\usepackage{amsmath,amssymb,mathtools,unicode-math}
\setmathfont{Latin Modern Math}
\usepackage{enumitem,needspace,xurl,hyperref}
\setlength{\parindent}{0pt}
\setlength{\parskip}{5pt}
\setlength{\emergencystretch}{3em}
\newcommand{\N}{\mathbb{N}}
\newcommand{\Z}{\mathbb{Z}}
\newcommand{\Q}{\mathbb{Q}}
\newcommand{\R}{\mathbb{R}}
\newcommand{\C}{\mathbb{C}}
\newcommand{\F}{\mathbb{F}}
\newcommand{\A}{\mathbb{A}}
\newcommand{\PP}{\mathbb{P}}
\newcommand{\E}{\mathbb{E}}
\newcommand{\eps}{\varepsilon}
\newcommand{\dd}{\,\mathrm{d}}
\begin{document}
\section*{UnsolvedMath ID: 7200057}
\textbf{Source problem code:} catalogue entry
\subsection{Definitions and mathematical statement}
In the usual complete-graph interpretation, let $S\subset\R^2$ be an $n$-point set in general position, with no three points collinear. Draw every edge of $K_n$ as its straight segment. Let $c(S)$ count unordered pairs of edges crossing in their relative interiors, and define
\[
\overline{\operatorname{cr}}(K_n)=\min_{|T|=n}c(T),
\]
where the minimum ranges over such planar point sets. Crossing pairs, rather than merely distinct intersection locations, are counted.

For even $n=2m$, a halving line is a line through two points of $S$ leaving exactly $m-1$ other points in each open half-plane. Count each unoriented line once, let $h(S)$ be their number, and put $h_{\max}(n)=\max_{|T|=n}h(T)$ in the same admissible class. The source's extremal question then becomes
\[
\forall n\text{ even}\ \forall S,\qquad
c(S)=\overline{\operatorname{cr}}(K_n)\quad\Longrightarrow\quad h(S)=h_{\max}(n).
\]
Thus it asks whether every crossing-minimizing configuration maximizes the number of halving lines, not merely whether some configuration simultaneously optimizes both quantities.

The source does not explicitly specify $K_n$, general position, parity, or the halving-line convention. For odd $n$, equal numbers on the two sides of a point-pair line are impossible; almost-halving pair lines and bisectors through a single point are different extensions. The display is therefore the standard even-cardinality candidate, and an all-$n$ version requires an explicit odd-cardinality definition.
\subsection{Short English statement}
Does every straight-line drawing of $K_n$ with the fewest crossings also have the maximum possible number of halving lines? The standard pair-line formulation uses even-sized point sets in general position.
\subsection{Sources}
\begin{itemize}
\item Local UnsolvedMath registry, permanent ID 7200057, exact source code catalogue entry; no proposer or year supplied. Registry status open, classification OPEN-TRIAGE.
\item Wikipedia contributors, \emph{List of unsolved problems in mathematics}, Geometry bullet 57, pinned revision 1366636767, \url{https://en.wikipedia.org/w/index.php?title=List_of_unsolved_problems_in_mathematics&oldid=1366636767#Geometry}. Supplied MediaWiki extraction accessed July 29, 2026; no primary attribution for the exact extremal implication is given.
\item The supplied summary cites ``On $(\leq k)$-edges, crossings, and halving lines of geometric drawings of $K_n$'', \texttt{arXiv:1102.5065}, for related edge/crossing identities, with no author list. It reports only abstract-level inspection and no verified 2024--2026 resolution of the exact optimizer implication; related identities are not treated as a proof.
\item Provenance flags primary-source status review NEEDS\_REVIEW and says the bullet remains unpublished pending that review, despite published=true metadata. Wikipedia text is CC BY-SA 4.0; attribution and share-alike handling remain required for release.
\end{itemize}
\end{document}
